\documentclass[11pt]{article}
\usepackage[utf8]{inputenc}
\usepackage[T1]{fontenc}
\usepackage{graphicx}
\usepackage{xcolor}
\usepackage{amsmath, amssymb}
\usepackage{url}
\usepackage{natbib}
\setlength{\parindent}{0em}
\usepackage[letterpaper, total={6.5in, 9in}]{geometry}
\newcommand*{\mybox}[2]{\colorbox{#1!30}{\parbox{.98\linewidth}{#2}}}
\setlength{\parskip}{0.5em}

\usepackage{setspace}
\onehalfspacing

\usepackage{xr-hyper}   
\makeatletter
\newcommand{\myexternaldocument}[2][]{%
  \externaldocument[#1]{#2}%
  \IfFileExists{#2.aux}{}{%
    \typeout{No file #2.aux.}%  
  }%
}
\IfFileExists{newcommands.tex}{% --- LaTeX Packages ---
% --- LaTeX Packages ---
\usepackage{amsmath}
\usepackage{amsthm}
\usepackage{bm}            % For bold math symbols
\usepackage{graphicx}      % For including images
\usepackage{xcolor}        % Required for \highlight
\usepackage{xspace}        % Required for commands like \nref
\usepackage{tikz}
%\usetikzlibrary{arrows.meta}
%\usepackage{mathtools}
%\usepackage{filecontents}
\usepackage{booktabs}
%\usepackage{placeins}
\usepackage{hyperref}
        \hypersetup{
          colorlinks=true,
          linkcolor=blue,
          urlcolor=blue,
          citecolor=blue
        }
\usepackage{etoolbox}

\setlength{\bibsep}{5pt plus 5pt} % space BETWEEN items
% For tighter lines WITHIN each item:
\usepackage{setspace}
\AtBeginEnvironment{thebibliography}{\setstretch{1.05}}

% --- Custom Commands ---

% Acronyms and Custom Text
\newcommand{\CVIC}{\text{CVIC}}
\newcommand{\deviance}{\text{dev}}
\newcommand{\dpois}{\text{dpois}}
\newcommand{\ic}{\text{ic}}
\newcommand{\nmsp}{\text{\scriptsize NMSP}}
\newcommand{\nrsp}{\text{\scriptsize NRSP}}
\newcommand{\pmin}{p_{\text{\scriptsize min}}}
\newcommand{\pp}{\text{pp}}
\newcommand{\nref}{[\textbf{references}]}
\newcommand{\SP}{survival probability\xspace}
\newcommand{\SPs}{survival probabilities\xspace}
\newcommand{\svf}{S_{i}(\cdot)}
\newcommand{\unif}[0]{\text{Unif}}
\newcommand{\rsp}[0]{\text{RSP}}
\newcommand{\msp}[0]{\text{MSP}}
\newcommand{\PPP}[0]{\text{PPP}}
\newcommand{\PPPs}[0]{\text{PPPs}}
\newcommand{\PPC}[0]{\text{PPC}}
\newcommand{\bU}[0]{\bm{U}}
\newcommand{\pv}[0]{\mathcal{P}}

% Math Symbols (Bold)
\newcommand{\btheta}[0]{\bm{\theta}}
\newcommand{\bTheta}[0]{\bm{\Theta}}
\newcommand{\bx}[0]{\bm{x}}
\newcommand{\by}[0]{\bm{y}}
\newcommand{\bY}[0]{\bm{Y}}
\newcommand{\bbeta}{\bm{\beta}}
\newcommand{\bBeta}{\bm{\Beta}}
\newcommand{\blambda}{\bm{\lambda}}
\newcommand{\bmu}{\bm{\mu}}
\newcommand{\bphi}{\bm{\phi}}
\newcommand{\bp}{\bm{p}}
\newcommand{\bsgm}{\bm{\sigma}}
\newcommand{\hu}[0]{\pi^o_{i}}
\newcommand{\mhu}[0]{\pi^o}
% General Superscripts and Subscripts
\newcommand{\obs}[0]{^{\text{obs}}}
\newcommand{\sobs}[0]{^{\text{obs}}}
\newcommand{\post}[0]{^{\text{Post}}}
\newcommand{\cv}[0]{^{\text{CV}}}
\newcommand{\iscv}[0]{^{\text{ISCV}}}
\newcommand{\supt}[0]{^{(t)}}
\newcommand{\mci}{^{(t)}}        % Kept as it has a different name from \supt
\newcommand{\mcs}{^{(s)}}
\newcommand{\rep}{^{\text{rep}}}
\newcommand{\ppc}{_{\text{ppc}}}
\newcommand{\noi}{_{-i}}
\newcommand{\wi}{_{1:n}}
\newcommand{\IS}{^{\text{\scriptsize IS}}}

% Compound Superscripts for Z-residuals
\newcommand{\rspp}[0]{^{\text{Post-RSP}}}
\newcommand{\mspp}[0]{^{\text{Post-MSP}}}
\newcommand{\rspis}[0]{^{\text{ISCV-RSP}}}
\newcommand{\mspis}[0]{^{\text{ISCV-MSP}}}
\newcommand{\rsppost}{^{\text{Post-RSP }}}  % From first block
\newcommand{\rspiscv}{^{\text{ISCV-RSP}}}    % From first block
\newcommand{\E}{\mathbb{E}}
\newcommand{\Epost}{\mathbb{E}_{\btheta\,|\,\by}}
\newcommand{\Ecv}{\mathbb{E}_{\btheta\,|\,\by_{-i}}}
\newcommand{\highlight}[1]{\textcolor{blue}{#1}}
\newcommand{\pvalue}[0]{$p$-value\xspace}
\newcommand{\pvalues}[0]{$p$-values\xspace}
\newcommand{\norm}[1]{\left\lVert#1\right\rVert}

\newcommand{\textred}[1]{\textcolor{orange}{#1}}


% ===================================================================
%                 THEOREM AND PROOF DEFINITIONS
% ===================================================================



%old theorem style for ba

% % --- 1. Define Custom Theorem Styles ---
% % A style for theorems where the optional note is bold
% \newtheoremstyle{boldnote}
%   {}		    % Space above
%   {}		    % Space below
%   {\itshape}	% Body font
%   {}		    % Indent amount
%   {\bfseries}	% Theorem head font
%   {.}		    % Punctuation after head
%   {.5em}	    % Space after head
%   {\thmname{#1}\thmnumber{ #2}\thmnote{ [\textbf{#3}]}}


% % --- 2. Define Theorem-Like Environments ---
% % Set the default style for all environments defined from this point
% \theoremstyle{plain}

% % The main theorem environment, numbered per section. All others will follow its counter.
% \newtheorem{theorem}{Theorem}[section] 

% % Environments that share the theorem counter
% \newtheorem{lemma}[theorem]{Lemma}
% \newtheorem{corollary}[theorem]{Corollary}

% % Switch to the custom style for any environments you want to use it
% % \theoremstyle{boldnote}
% % \newtheorem{specialtheorem}[theorem]{Special Theorem} % Example


% % --- 3. Define Proof and Step Environments ---
% % For custom "Step" counter within proofs
% \newtheorem{proofpart}{Step}

% % Change the name "Proof" to "Proof" (bold)
% \renewcommand{\proofname}{\textbf{Proof}}

% % Automatically reset the 'Step' counter at the beginning of every proof
% \AtBeginEnvironment{proof}{\setcounter{proofpart}{0}}
% % ===================================================================
}{}

\makeatother


\myexternaldocument[S-]{supplement}


\makeatother






\begin{document}


\begin{center}
    {\LARGE \textbf{Response to Reviewer A}} \\[1em]
    {\large \today}
\end{center}
\vspace{2em} % Space between title and your text
\normalsize
\normalsize

We sincerely thank you for your time, valuable feedback, and general appreciation of the paper. Guided by the collective insights from you and other referees, we have substantially revised the manuscript. We have responded to your comments item by item, and your constructive suggestions have greatly improved both the methodology and the presentation. In this detailed response, we quote several updated passages for your reference (please note that some of these quoted excerpts may have been further refined in the final manuscript to ensure perfect flow and accuracy). In the revised manuscript, revised text is highlighted in \textred{orange}; for substantially rewritten sections, only the \textred{section title} is highlighted rather than the entire text. 

\section*{General Comments}

\mybox{gray}{\textbf{General comment.} This is an interesting piece of work. The authors propose a $Z$-residual concept that applies to general Bayesian models. The following issues need to be addressed. The major issues are \#3, \#4, \#6, and \#7.}

\medskip

\mybox{gray}{1. Page 3, line 36. Can the authors explain why you use the term ``$Z$-residual''?}

\paragraph{Response:}
Thank you for this helpful suggestion. We use the term “Z-residual” because the proposed residual is represented on the standard-normal, or Z-score, scale. Specifically, each observation is first transformed to a predictive probability scale through the randomized survival probability, and this probability is then mapped to the standard-normal scale by the normal-score transformation, ($z_i=-\Phi^{-1}{RSP_i(y_i)}$).
Under correct model specification, the parameter-wise randomized survival probability is uniform, so the corresponding residual has a standard normal reference distribution. Thus, the term “Z-residual” emphasizes the normal-score scale of the residual, rather than implying that it is a classical additive residual. 

 We have clarified this naming in Section 2.2, in the paragraph immediately following equation (2), by adding a sentence explaining that the term “Z-residual” refers to this standard-normal reference scale.


\medskip

\mybox{gray}{2. Page 3, line 40. Can the authors clarify if $U_i$ is needed in the randomized survival probability when the outcome variable $Y$ is continuous?}

\paragraph{Response:}
Thank you for this helpful comment. When (Y) is continuous, the auxiliary random variable ($U_i$) is not needed. The randomization term is introduced to handle probability mass at observed values in discrete outcomes. For a continuous outcome, the probability mass at any exact observed value is zero, so ($p_i(y_i\mid\theta)=0$) and the randomized survival probability reduces to the usual survival-probability transform. Equivalently, the construction becomes a probability-integral-transform type residual followed by the normal-score transformation. 

 We have clarified this point in Section 2.2, in the paragraph immediately following equation (1), by stating that ($U_i$) is only needed when the Response:has probability mass at the observed value and that, for continuous Response:, the construction reduces to ($S_i(y_i\mid\theta)$). We also reiterate this point in the first paragraph of Section 3.2 when introducing the continuous inverse Gaussian simulation.

\medskip

\mybox{gray}{3. \textbf{(How to position this paper?)} If $U_i$ is not needed when $Y$ is continuous, the authors need to clarify the main contribution of the paper. When $Y$ is continuous, is your residual still novel in the Bayesian literature? Does the significance still remain? If not, the authors need to reposition this paper by saying that the novelty resides in the discrete case.}

\paragraph{Response:}

Thank you for raising this important point. We agree that for continuous responses, the auxiliary randomization $U_i$ is unnecessary, and the pointwise probability transform is simply the probability integral transform. 

We entirely agree that the PIT and LOO-PIT are not novel concepts, and while Z-residuals (randomized quantile residuals) have been recently discussed much in frequentist statistics in the past decade, they remain largely absent from Bayesian model assessment. Therefore, our primary methodological contribution lies not in inventing these foundational concepts, but in synthesizing them into a framework for constructing \textbf{observation-level residual diagnostics} in applied Bayesian modeling. By combining randomized residuals with observation-level summaries over posterior or LOO predictive distributions, we address a critical, overlooked gap: the shift from overall model \textit{checking} to granular residual \textit{diagnostics}. Traditional goodness-of-fit (GOF) tests typically evaluate the overall distribution or rely on a single aggregate $p$-value to merely confirm or reject a model. In contrast, our expanded framework---supported by theoretical justification for its asymptotic calibration---facilitates rich exploratory analyses using standard graphical tools (e.g., scatterplots, boxplots, and quantile-quantile plots) alongside formal statistical tests. This enables practitioners to systematically pinpoint and understand specific structural flaws within the model. These tests include, but are not limited to: the ANOVA test for misspecified covariate functional forms (introduced in this paper); the Box-Ljung test for temporal autocorrelation; Moran's $I$ statistic for global spatial dependencies; and the Breusch-Pagan and Bartlett tests for heteroskedasticity. Ultimately, the proposed Z-residuals enable the seamless incorporation of the full suite of frequentist diagnostic tools into Bayesian analysis.

To bring the rich suite of residual diagnostic tools into Bayesian statistics, we introduce the following \textbf{specific methodological innovations} in this paper:
\begin{itemize}
    \item \textbf{A ``summarize-first, then-test'' paradigm:} We construct Z-residuals by integrating randomized survival probabilities (RSPs) over the posterior distribution prior to applying conventional diagnostic tools. This fundamentally replaces the traditional posterior predictive check (PPC) ``test-first, then-summarize'' approach, which evaluates discrepancy statistics for each MCMC draw and then summarizes the resulting $p$-values. Despite its conceptual simplicity, this reversal yields a diagnostic tool with solid theoretical foundations that unifies frequentist and Bayesian workflows. Notably, the construction of Z-residuals is independent of any specific test; therefore, they can be subsequently analyzed using any graphical or numerical diagnostic tool.
    \item \textbf{Frequentist asymptotic justification:} We provide a rigorous theoretical proof demonstrating that the posterior and LOOCV RSPs are asymptotically uniformly distributed, and the resulting Z-residuals are asymptotically normally distributed from a strictly frequentist perspective. This theoretical grounding validates the application of existing frequentist residual diagnostic tools to Bayesian models. 
    \item \textbf{A distribution-free upper-bound $p$-value:} To overcome the auxiliary randomization introduced for discrete responses, we provide a novel, distribution-free upper bound to reliably summarize replicated $p$-values (or, more broadly, dependent $p$-values). This technique has broad applicability beyond the current context.
\end{itemize}

\textbf{In terms of empirical findings}, this paper establishes the following critical points: 

\begin{itemize}
    \item Relying solely on the marginal distribution of PIT or Z-residuals for overall goodness-of-fit (GOF) checking is fundamentally inadequate. As we demonstrate, covariate-specific tests---such as the Z-AOV---successfully detect functional form misspecifications, highlighting how Z-residuals can be leveraged to diagnose a wide array of structural flaws. These subtle deficiencies are often completely masked when inspecting only the marginal distribution of PITs or relying on simple summary statistics that ignore relationships with covariates.
    
    \item Graphical diagnostics utilizing Z-residuals serve as powerful exploratory tools for \textit{discovering} unknown model deficiencies. Crucially, these visual insights guide the subsequent application of formal statistical tests, whether employing posterior/ISCV Z-residuals or conducting Z-residuals-PPC tests.
    
    \item Tests based on Z-residuals are well-calibrated and yield significantly higher statistical power than both traditional PPCs utilizing the $\chi^2$ statistic and Z-residual-based PPCs. This dual advantage in calibration and power is attributed to two key factors: (1) unlike arbitrary discrepancy measures, the exact distribution of Z-residuals is known under the true predictive distribution; and (2) our ``summarize-first, then-test'' paradigm asymptotically approximates this true predictive distribution, as formally justified by the theory presented in this paper.
    
    \item Our simulations demonstrate that  Z-residual-based PPCs mitigate the notorious conservatism of traditional $\chi^2$ PPCs. Because they average over the MCMC draws, they remain highly stable despite the auxiliary $\bm{U}$ drawn at each iteration, serving as a reliable secondary confirmation of model misfit.
\end{itemize}

 We have substantially revised Section 1 and Section 5, incorporating the points above to more clearly articulate our contributions and findings.

\medskip

\mybox{gray}{4. \textbf{(Lack of sensitivity analysis)} Does the choice of the prior influence diagnosis? In particular, when the sample size is small, an ill-chosen prior may lead to a wrong rejection of the model. Am I right? The authors should perform sensitivity analysis.}

\paragraph{Response:}

Thank you for raising this important point. You are absolutely correct that if a very poor or ill-chosen prior is used, particularly with a small sample size, it can lead to the rejection of the model. However, from a Bayesian perspective, the prior is an integral component of the model specification. If an inappropriate prior severely degrades the model's predictive fit, rejecting the model is not a "wrong" rejection, but exactly a type of misspecification our diagnostic tools are devised to identify. Furthermore, because Z-residuals focus specifically on the predictive distribution of the observed data $y_i$, they are generally robust to routine prior choices. The diagnostic is not overly sensitive to prior misspecification unless the prior is highly influential and strictly conflicts with the likelihood. In most practical scenarios, the data will adequately inform the posterior predictive distribution, keeping the Z-residuals stable.


To illustrate the robustness of our diagnostic conclusions to the choice of prior, we conducted a sensitivity analysis on the dolphin-sighting data under baseline, tighter, and more diffuse prior specifications (Supplementary Table~S5). After refitting the models and recomputing all diagnostic $p$-values (Z-residuals, PPC, and HPC), we found the qualitative patterns to be highly stable (Supplementary Figure~S4). Most notably, the diagnostics consistently reject the linear HNB and NB models while accepting their nonlinear counterparts across all three prior regimes. We have added this full analysis to the Supplementary Material and noted the robustness of these findings in Section~4.

\medskip

\mybox{gray}{5. Page 17, Figure 1. I am not sure about the value of plotting $Z$-residuals against the index. What insights can the plots offer in addition to those offered by the QQ plots?}

\paragraph{Response.}
Thank you for this helpful comment. We agree with your assessment and have removed the residual-versus-index plots, replacing them with residual-versus-covariate plots in the revised manuscript. While index plots can sometimes be useful in practice—since data indices are often non-random and may reveal order-dependent patterns or temporal clustering—we completely agree that for our specific simulation studies, the index contains no meaningful information.


\medskip

\mybox{gray}{6. \textbf{(Lack of numerical studies with respect to other Bayesian models)} The authors claim that they proposed ``general-purpose residuals''. But the numerical studies only focus on (hurdle) negative binomial models. This is relevant to my comment \#3. If the authors position this paper as a general methodology, then more numerical studies using different models (for continuous and discrete data) should be presented.}

\paragraph{Response.}

Thank you for this important suggestion. We agree that focusing solely on discrete count models limited our broader methodological claims. To address this and provide more diverse experimental scenarios, we have substantially broadened our numerical studies:
\begin{itemize}
\item \textbf{Added a Continuous Inverse Gaussian (IG) Simulation (Section~3.2):} We replaced the original count-data simulation in the main text with a continuous IG regression study. This example explicitly demonstrates the fundamental performance differences between PPCs and Z-residuals without the potential interference of the auxiliary randomization required for discrete outcomes.
\item \textbf{Added a Mixed Misspecification Simulation (Section~S1.2):} Addressing your advice for broader experimental scenarios, and motivated by another reviewer's feedback, we added a new simulation to the Supplementary Material that explores simultaneous misspecifications in both distributional families and covariate functional forms.
\item \textbf{Relocated the Original Simulation (Section~S1.1):} The original HNB-versus-NB wrong-family simulation has been moved to the Supplementary Material to ensure the main text focuses on the broader applicability of the methods.
\end{itemize}


\medskip

\mybox{gray}{7. \textbf{(Double use of the data)} I like the point the authors make about the double use of data in fitting and diagnosing the model. My question is: although you do not use $y_i$, you do use the rest of the data. Why does this not count as ``double use of the data''? Furthermore, $y_i$ is generated from the same model that generates other $y$'s. So I believe your approach still has the same problem. The importance-sampling approach presented in Section 2.4 may alleviate the severity of the problem.}

\paragraph{Response.}
Thank you for this insightful comment. You are absolutely correct: both the proposed posterior and cross-validatory (ISCV) Z-residuals technically reuse the data. The fundamental distinction between our approach and standard Posterior Predictive Checks (PPCs) lies not in the complete elimination of data reuse, but in the \textit{severity} of the bias it causes, particularly in large samples.

To clarify this distinction, we have revised Section 2.7 to summarize the asymptotic differences:

\begin{quote}
    
\begin{itemize}
    \item \textbf{Tests on Z-residuals:} Z-residuals are constructed observation-by-observation. As the sample size $n\to \infty$, the influence of any single observation $y_i$ on the full posterior predictive distribution vanishes. Consequently, the full-data posterior predictive law and the LOO predictive law become asymptotically indistinguishable, both converging to the true data-generating law. Thus, the posterior and cross-validatory Z-residuals are asymptotically independent and follow a $N(0,1)$ distribution.
    
    \item \textbf{PPCs:} In contrast, a PPC evaluates a dataset-level discrepancy statistic $T(\mathbf{y}, \boldsymbol{\theta})$. Because the \textit{entire} dataset $\mathbf{y}$ is used simultaneously to simulate the reference distribution, the simulated data is systematically pulled toward the observed data. This dataset-level double use causes PPC $p$-values to be inherently conservative even asymptotically \citep{Robins2000}. Specifically, as the sample size $n \to \infty$, the downward bias in the expected discrepancy does not wash out; rather, it converges to an $O(1)$ constant determined by the parameter count, $p$. Conceptually, this phenomenon parallels the optimism bias formalized by Wilks' theorem in frequentist statistics: just as maximum likelihood estimation absorbs data noise to produce a non-vanishing $p$-dimensional upward bias in the log-likelihood, the double use of data in PPCs systematically shrinks the expected discrepancy by $p$ degrees of freedom. Because this penalty is permanent, it severely limits the statistical power of standard PPCs to detect subtle misfits.
\end{itemize}
\end{quote}

Regarding your specific point about the cross-validatory approach: we agree that because $y_i$ and $y_j$ are generated from the same model, and their respective training sets ($y_{-i}$ and $y_{-j}$) share $n-2$ observations, this overlapping data inherently introduces finite-sample correlation between the LOO Z-residuals. We have revised the text to explicitly state that ISCV alleviates the direct self-influence of $y_i$, but does not eliminate this correlation.

However, despite these finite-sample dependencies, our simulation studies demonstrate that the practical impact of data double-use in posterior Z-residuals is remarkably minor when aggregate diagnostic tests (like Shapiro--Wilk and ANOVA) are employed. The $p$-values from these tests remain nearly uniform under the null, showing vastly superior calibration and power compared to standard PPCs (especially those relying on the common $\chi^2$ statistic). 

Finally, we have added a note clarifying that when data is highly abundant, strict data-splitting to reserve a hold-out test set remains the most definitive way to completely eradicate data double-use.

\medskip

\mybox{gray}{8. \textbf{(Insufficient simulation replications)} Page 20, Figure 3. The authors should increase the simulation replications to 1000 if type I errors are part of the results.}

\paragraph{Response.}
Thank you for this suggestion. While 1000 replications is ideal, the massive computational burden of Bayesian modeling---where a single MCMC fit can take 20 mins---makes this prohibitive across our expanded simulation suite. To practically estimate the rejection rates, we increased all scenarios to 500 replications. At $R=500$, the Monte Carlo standard error for a nominal 0.05 Type I error rate is approximately $\sqrt{0.05(1-0.05)/500} \approx 0.010$, yielding a 95\% margin of error of roughly $\pm 0.019$. For a rejection rate near 0.50 (where variance is maximized), the standard error is $\sqrt{0.50(1-0.50)/500} \approx 0.022$, corresponding to a margin of error of about $\pm 0.044$. This precision is sufficient for discerning diagnostic behaviors across different methods, and we have detailed these uncertainty calculations in the Supplementary Material.

\medskip

\mybox{gray}{9. \textbf{(Inflated type I error)} Page 20, Figure 3(a). Please report type I errors in numerical form. By reading from the graphical presentation, I feel that some of the type I errors are inflated, close to 0.1. The same can be said for Figure 6.}

\paragraph{Response.}
Thank you for this valuable suggestion. We agree that graphical presentations alone are insufficient for accurately assessing Type I error inflation. We have fully taken your advice and compiled the exact numerical rejection rates for all simulation settings. Due to page limitations in the main manuscript, these detailed numerical summaries have been placed in the Supplementary Material. Please refer to Tables S1, S3 and S4, which explicitly report the empirical Type I error rates and statistical power, complete with Monte Carlo standard error calculations to clarify the Monte Carlo estimating errors.

\medskip

\mybox{gray}{10. \textbf{(Bin size in Figure 4)} How do the authors determine the bin size? The analysis of variance (and its $p$-value) is heavily dependent on how you bin the residual points.}

\paragraph{Response.}Thank you for pointing this out. We agree that the grouped ANOVA diagnostic depends on the binning rule, which must be made explicit. In the revised Section~2.5, we specify that for continuous covariates, we divide the observed range into $k$ equally spaced intervals, recommending $k=10$ as a moderate default to balance resolution against group-size stability.

To address your concern, we added a binning sensitivity analysis to the Supplementary Material (Table~S7, Figure~S6). Using the dolphin-sighting data, we evaluated $k \in \{5, 10, 15, 20\}$, summarizing 500 randomizations per $k$ via the median upper-bound $p$-value. The qualitative conclusions remain remarkably stable: the linear HNB model consistently shows residual patterns against depth, while the nonlinear HNB model does not.

This stability likely stems from the automatic degrees-of-freedom (df) adjustment inherent in the ANOVA test. When $k$ is small, the captured difference in group means is smaller, but the df for the Mean Squared Error (MSE) is larger. Conversely, when $k$ is large, the mean differences may be more pronounced, but the MSE df decreases. In practice, we recommend fixing $k=10$ as a default, though checking a few different values of $k$ remains a sensible strategy to verify whether a functional form misspecification is present.


\medskip

\mybox{gray}{11. \textbf{(Tone down the self-promoting statements)} Page 15, line 40: ``very promising''. Page 30, line 26: ``excellent performance''. Please use objective statements.}

\paragraph{Response.}
Thank you for this valuable suggestion. We agree that certain expressions in the original manuscript were overly promotional. Throughout the revised manuscript---particularly in the Abstract and Sections 1, 3, 4, and 5---we have replaced subjective phrases (e.g., ``very promising,'' ``excellent performance,'' ``superior'') with neutral, objective alternatives (e.g., ``useful,'' ``favorable,'' ``better in the settings considered''). Furthermore, we carefully revised all comparative statements to ensure they are strictly grounded in our specific simulation or data-analysis settings, rather than reading as broad, generalized claims.
\medskip

\mybox{gray}{12. \textbf{(Reduce the excessive use of uncommon abbreviations)} The authors used GOF, PPC, PPP, RSP, LOO, ISCV, NB, HNB, MSP, SW, PIT, AOV, and other uncommon abbreviations. They do make the reading difficult. I have to flip the pages and spend a lot of time to find what they really mean. The journal \textit{Biometrika} explicitly forbids the use of uncommon abbreviations. Please remove as many as possible.}

\paragraph{Response.}


Thank you for this helpful comment. We agree that excessive abbreviations hindered readability, and we have significantly reduced their use throughout the revised manuscript. We now write out low-frequency abbreviations (e.g., CDF, PMF, LOOCV) and general terms (e.g., goodness-of-fit) in full. However, to prevent excessive length and maintain clean layouts, we retained a small set of core, high-frequency abbreviations (e.g., RSP, PPC, ISCV, HNB, SW). These are clearly defined at first use and, for the reader's convenience, \textit{we have now added a list of these core abbreviations to the first page of the manuscript.} Finally, we updated figure and table captions to explicitly spell out key tests (such as the Shapiro--Wilk and ANOVA tests) so they function as self-contained references without requiring readers to search for definitions.

\setlength{\bibsep}{0.2em}
\bibliographystyle{agsm}
\bibliography{ref}

\end{document}